% Ban thu: chi giu cac cau co trong CD2

%@ id=TT25.04.003 | trang=5 | so=2 | de=II | vitri=CD2.4,C3.9.5 | muc=8 | tt=0
\begin{bt}
Cho hai biểu thức $A = \dfrac{x+2\sqrt{x}}{x-1}$ và $B = \dfrac{2}{x} - \dfrac{2-x}{x(\sqrt{x}+1)}$ với $x > 0; x \ne 1$.
\begin{cacy}
\item Tính giá trị của biểu thức $A$ khi $x = 25$.
\item Chứng minh: $B = \dfrac{\sqrt{x}+2}{\sqrt{x}(\sqrt{x}+1)}$.
\item Tìm $x$ để biểu thức $\dfrac{A}{B} > 1$.
\end{cacy}
\loigiai{
\begin{cacy}
\item Thay $x = 25$ (thỏa mãn ĐK) vào biểu thức $A$ ta được
\[ A = \dfrac{25+2\sqrt{25}}{25-1} = \dfrac{25+2 \cdot 5}{24} = \dfrac{35}{24} \]
Vậy khi $x = 25$ thì $A = \dfrac{35}{24}$.
\item Điều kiện: $x > 0; x \ne 1$.
\[ B = \dfrac{2}{x} - \dfrac{2-x}{x(\sqrt{x}+1)} = \dfrac{2\sqrt{x}+2-2+x}{x(\sqrt{x}+1)} \]
\[ B = \dfrac{x+2\sqrt{x}}{x(\sqrt{x}+1)} = \dfrac{\sqrt{x}(\sqrt{x}+2)}{x(\sqrt{x}+1)} = \dfrac{\sqrt{x}+2}{\sqrt{x}(\sqrt{x}+1)} \]
\item Ta có $\dfrac{A}{B} > 1 \Rightarrow \dfrac{x+2\sqrt{x}}{x-1} : \dfrac{\sqrt{x}+2}{\sqrt{x}(\sqrt{x}+1)} > 1 \Rightarrow \dfrac{\sqrt{x}(\sqrt{x}+2)}{x-1} \cdot \dfrac{\sqrt{x}(\sqrt{x}+1)}{\sqrt{x}+2} > 1$
\[ \Rightarrow \dfrac{x}{\sqrt{x}-1} > 1 \Rightarrow \dfrac{x-\sqrt{x}+1}{\sqrt{x}-1} > 0 \quad (*) \]
Xét $x - \sqrt{x} + 1 = \left(\sqrt{x}-\dfrac{1}{2}\right)^2 + \dfrac{3}{4} > 0$ với mọi $x > 0; x \ne 1$.

Do đó $(*) \Rightarrow \sqrt{x}-1 > 0 \Rightarrow x > 1$.

Kết hợp điều kiện, ta có $x > 1$ thỏa mãn yêu cầu bài toán.
\end{cacy}
}
\end{bt}
